\paragraph{规范体积密度、纤维奇偶荷与可解性阈值}
在定理 \ref{thm:xi-foldbin-collision-probability-normalized-groupoid-dimension} 的碰撞公式与定理 \ref{thm:xi-gauge-volume-boltzmann-stirling-exponential} 的指数律之间，还存在一条由 escort 相对熵控制的精确桥接；同时，定义 \ref{def:fold-bin-gauge-sign} 的纤维奇偶荷可把 Abel 化、中心、window-\(6\) 的全自同构结构以及低阶上同调统一为纯纤维判别。

\begin{theorem}[规范体积密度的二次 escort--KL 精确分解]\label{thm:xi-foldbin-gauge-escort-kl-decomposition}
\leanverified{paper\_xi\_foldbin\_gauge\_escort\_kl\_decomposition}
沿用命题 \ref{prop:fold-bin-chi2-col} 与命题 \ref{prop:fold-bin-gauge-volume} 的记号。记
\[
p_m^{\mathrm{bin}}(w):=\frac{d_m^{\mathrm{bin}}(w)}{2^m},
\qquad
\widetilde p_m(w):=\frac{\bigl(p_m^{\mathrm{bin}}(w)\bigr)^2}{\mathsf{Col}^{\mathrm{bin}}_m}
\qquad (w\in X_m).
\]
则有严格恒等式
\[
\boxed{\;
\kappa_m
=
\log\!\bigl(2^m\mathsf{Col}^{\mathrm{bin}}_m\bigr)
-D_{\mathrm{KL}}\!\bigl(p_m^{\mathrm{bin}}\|\widetilde p_m\bigr).
\;}
\]
等价地，
\[
\boxed{\;
\kappa_m
=
\log\!\Bigl(\frac{2^m}{|X_m|}\Bigr)
+\log\!\Bigl(1+\chi^2\!\bigl(p_m^{\mathrm{bin}}\|u_m\bigr)\Bigr)
-D_{\mathrm{KL}}\!\bigl(p_m^{\mathrm{bin}}\|\widetilde p_m\bigr),
\;}
\]
其中 \(u_m(w)=1/|X_m|\)。
特别地，
\[
\log\!\Bigl(\frac{2^m}{|X_m|}\Bigr)
\le
\kappa_m
\le
\log\!\bigl(2^m\mathsf{Col}^{\mathrm{bin}}_m\bigr),
\]
且两端取等都当且仅当 \(p_m^{\mathrm{bin}}=u_m\)。
\end{theorem}
\begin{proof}
为简记起见，写
\[
d_w:=d_m^{\mathrm{bin}}(w),
\qquad
p_w:=p_m^{\mathrm{bin}}(w),
\qquad
H(p):=-\sum_{w\in X_m}p_w\log p_w.
\]
由于 \(X_m\) 由实际出现的稳定类型组成，故对每个 \(w\in X_m\) 都有 \(d_w\ge 1\)，从而 \(p_w>0\)。
由定义
\[
\kappa_m
=
2^{-m}\sum_{w\in X_m}d_w\log d_w
=
\sum_{w\in X_m}p_w\log(2^m p_w)
=
m\log 2-H(p).
\]
另一方面，
\[
\begin{aligned}
D_{\mathrm{KL}}(p\|\widetilde p)
&=
\sum_{w\in X_m}p_w\log\frac{p_w}{p_w^2/\mathsf{Col}^{\mathrm{bin}}_m}\\
&=
\log\!\bigl(\mathsf{Col}^{\mathrm{bin}}_m\bigr)-\sum_{w\in X_m}p_w\log p_w\\
&=
\log\!\bigl(\mathsf{Col}^{\mathrm{bin}}_m\bigr)+H(p).
\end{aligned}
\]
消去 \(H(p)\) 即得
\[
\kappa_m
=
\log\!\bigl(2^m\mathsf{Col}^{\mathrm{bin}}_m\bigr)
-D_{\mathrm{KL}}(p\|\widetilde p).
\]
再由命题 \ref{prop:fold-bin-chi2-col} 的恒等式
\[
\chi^2\!\bigl(p_m^{\mathrm{bin}}\|u_m\bigr)
=
|X_m|\,\mathsf{Col}^{\mathrm{bin}}_m-1,
\]
便得到第二种写法。

下界由 Shannon 熵上界 \(H(p)\le \log|X_m|\) 立即推出，即
\[
\kappa_m=m\log 2-H(p)\ge m\log 2-\log|X_m|.
\]
且等号成立当且仅当 \(p=u_m\)。
上界则由 \(D_{\mathrm{KL}}(p\|\widetilde p)\ge 0\) 得到；若上界取等，则 \(p=\widetilde p\)，即
\[
p_w=\frac{p_w^2}{\mathsf{Col}^{\mathrm{bin}}_m}
\qquad (w\in X_m).
\]
由于所有 \(p_w>0\)，故 \(p_w=\mathsf{Col}^{\mathrm{bin}}_m\) 对所有 \(w\) 同时成立，于是 \(p\) 在 \(X_m\) 上均匀，即 \(p=u_m\)。
反之若 \(p=u_m\)，则显然同时达到上下界。证毕。
\end{proof}

\begin{corollary}[规范体积密度的碰撞表达与常数主差额]\label{cor:xi-foldbin-gauge-collision-rigid-shift}
沿用定理 \ref{thm:xi-foldbin-gauge-escort-kl-decomposition} 的记号，则
\[
\boxed{\;
g_m
=
\log\!\bigl(2^m\mathsf{Col}^{\mathrm{bin}}_m\bigr)
-1
-D_{\mathrm{KL}}\!\bigl(p_m^{\mathrm{bin}}\|\widetilde p_m\bigr)
+O\!\Bigl(m\bigl(\tfrac{\varphi}{2}\bigr)^m\Bigr).
\;}
\]
等价地，
\[
\boxed{\;
g_m
=
\log\!\Bigl(\frac{2^m}{|X_m|}\Bigr)
+\log\!\Bigl(1+\chi^2\!\bigl(p_m^{\mathrm{bin}}\|u_m\bigr)\Bigr)
-1
-D_{\mathrm{KL}}\!\bigl(p_m^{\mathrm{bin}}\|\widetilde p_m\bigr)
+O\!\Bigl(m\bigl(\tfrac{\varphi}{2}\bigr)^m\Bigr).
\;}
\]
因此 \(g_m\) 与碰撞项 \(\log(2^m\mathsf{Col}^{\mathrm{bin}}_m)\) 的主差额由常数 \(1\) 固定，余项仅为指数小修正。
\end{corollary}
\begin{proof}
由定理 \ref{thm:fold-bin-gauge-volume-stirling}，
\[
g_m=\kappa_m-1+O\!\left(\frac{|X_m|\,m}{2^m}\right)
=
\kappa_m-1+O\!\Bigl(m\bigl(\tfrac{\varphi}{2}\bigr)^m\Bigr),
\]
其中用到了 \(|X_m|=F_{m+2}=\Theta(\varphi^m)\)。
将定理 \ref{thm:xi-foldbin-gauge-escort-kl-decomposition} 的两种表示分别代入即得。证毕。
\end{proof}

\begin{theorem}[纤维奇偶荷的分裂正合列与最小完备秩]\label{thm:xi-foldbin-parity-charge-split-exact-minrank}
沿用命题 \ref{prop:fold-bin-gauge-decomposition} 与定义 \ref{def:fold-bin-gauge-sign} 的记号。记
\[
G_m^{\mathrm{bin}}:=\mathcal{G}^{\mathrm{bin}}_m=\Aut(\Fold_m^{\mathrm{bin}}),
\qquad
d_w:=d_m^{\mathrm{bin}}(w),
\qquad
N_m:=\#\{w\in X_m:\ d_w\ge 2\}.
\]
将定义 \ref{def:fold-bin-gauge-sign} 中的
\(
\mathrm{sgn}_m:G_m^{\mathrm{bin}}\to (\ZZ/2\ZZ)^{X_m}
\)
与坐标投影
\[
(\ZZ/2\ZZ)^{X_m}
\longrightarrow
(\ZZ/2\ZZ)^{\{w\in X_m:\ d_w\ge 2\}}
\cong
(\ZZ/2\ZZ)^{N_m}
\]
复合后仍记为 \(\mathrm{sgn}_m\)。
约定 \(\mathfrak A_1=\mathfrak A_2=1\)。
则存在分裂正合列
\[
\boxed{\;
1\longrightarrow \prod_{w\in X_m}\mathfrak A_{d_w}
\longrightarrow G_m^{\mathrm{bin}}
\xrightarrow{\ \mathrm{sgn}_m\ }
(\ZZ/2\ZZ)^{N_m}
\longrightarrow 1.
\;}
\]
从而：
\begin{enumerate}
  \item \(\bigl(G_m^{\mathrm{bin}}\bigr)^{\mathrm{ab}}\cong (\ZZ/2\ZZ)^{N_m}\)；
  \item 若群同态 \(\Phi:G_m^{\mathrm{bin}}\to (\ZZ/2\ZZ)^r\) 满足存在线性映射
  \[
  L:(\ZZ/2\ZZ)^r\to (\ZZ/2\ZZ)^{N_m}
  \]
  使得 \(\mathrm{sgn}_m=L\circ\Phi\)，则必有 \(r\ge N_m\)；
  \item 上述下界由 \(\mathrm{sgn}_m\) 自身达到，因此 \(N_m\) 正是纤维奇偶荷的最小完备秩。
\end{enumerate}
特别地，在 \(m=6\) 时，由推论 \ref{cor:fold-bin-m6-fiber-hist} 的直方图 \(2:8,3:4,4:9\) 可得 \(N_6=21\)，于是
\[
\bigl(G_6^{\mathrm{bin}}\bigr)^{\mathrm{ab}}\cong (\ZZ/2\ZZ)^{21}.
\]
\end{theorem}
\begin{proof}
由命题 \ref{prop:fold-bin-gauge-decomposition}，
\[
G_m^{\mathrm{bin}}
\cong
\prod_{w\in X_m}\mathfrak S_{d_w}.
\]
于是只需逐纤维考察对称群。

当 \(d=1\) 时 \(\mathfrak S_1\) 平凡。
当 \(d=2\) 时 \(\mathfrak S_2\cong \ZZ/2\ZZ\)，其符号映射即为同构，核为平凡群，也即 \(\mathfrak A_2=1\)。
当 \(d\ge 3\) 时，有标准分裂正合列
\[
1\longrightarrow \mathfrak A_d
\longrightarrow \mathfrak S_d
\xrightarrow{\ \mathrm{sgn}\ }
\ZZ/2\ZZ
\longrightarrow 1,
\]
任取一换位即可给出截面。
对全部 \(w\in X_m\) 取直积，便得到所述分裂正合列。

第 (1) 条随即成立，因为右端已是所有因子 Abel 化的直积。
对第 (2) 条，条件 \(\mathrm{sgn}_m=L\circ\Phi\) 表明 \((\ZZ/2\ZZ)^{N_m}\) 是 \(\Phi(G_m^{\mathrm{bin}})\subseteq (\ZZ/2\ZZ)^r\) 的一个线性商。
由于 \(\mathrm{sgn}_m\) 满射，故
\[
\operatorname{rank}_{\FF_2}\Phi(G_m^{\mathrm{bin}})
\ge N_m,
\]
而 \(\Phi(G_m^{\mathrm{bin}})\) 是 \((\ZZ/2\ZZ)^r\) 的子空间，故其秩至多为 \(r\)，于是 \(r\ge N_m\)。
第 (3) 条由 \(\mathrm{sgn}_m\) 本身即为到 \((\ZZ/2\ZZ)^{N_m}\) 的满射立即得到。

最后，在 \(m=6\) 时，表 \ref{tab:fold_bin_degeneracy_spectrum_m6_m12} 的直方图给出全部 \(21\) 条纤维都满足 \(d_w\ge 2\)，故 \(N_6=21\)。证毕。
\end{proof}

\begin{theorem}[中心与可解性的纤维判别]\label{thm:xi-foldbin-center-solvable-fiber-criterion}
沿用命题 \ref{prop:fold-bin-gauge-decomposition} 的记号。记
\[
G_m^{\mathrm{bin}}:=\mathcal{G}^{\mathrm{bin}}_m,
\qquad
d_w:=d_m^{\mathrm{bin}}(w),
\qquad
B_m:=\#\{w\in X_m:\ d_w=2\},
\qquad
d_{\max}(m):=\max_{w\in X_m}d_w.
\]
则中心满足
\[
\boxed{\;
Z\!\bigl(G_m^{\mathrm{bin}}\bigr)\cong (\ZZ/2\ZZ)^{B_m}.
\;}
\]
此外，\(G_m^{\mathrm{bin}}\) 可解当且仅当 \(d_{\max}(m)\le 4\)。
若以 \(\ell_{\mathrm{der}}(G)\) 记群 \(G\) 的导出长度，则在可解情形下
\[
\ell_{\mathrm{der}}\!\bigl(G_m^{\mathrm{bin}}\bigr)
=
\begin{cases}
0,& d_{\max}(m)=1,\\
1,& d_{\max}(m)=2,\\
2,& d_{\max}(m)=3,\\
3,& d_{\max}(m)=4.
\end{cases}
\]
特别地，表 \ref{tab:fold_bin_degeneracy_spectrum_m6_m12} 的 \(m=6\) 行给出 \(B_6=8\) 与 \(d_{\max}(6)=4\)，故
\[
Z\!\bigl(G_6^{\mathrm{bin}}\bigr)\cong (\ZZ/2\ZZ)^8,
\qquad
\ell_{\mathrm{der}}\!\bigl(G_6^{\mathrm{bin}}\bigr)=3.
\]
并且定理 \ref{thm:xi-window6-center-exact-splitting-bdry-cyclic} 中的 boundary 中心子群
\[
(\ZZ/2\ZZ)^3_{\mathrm{bdry}}\hookrightarrow Z\!\bigl(G_6^{\mathrm{bin}}\bigr)
\]
正是该 \(8\) 维中心的规范直和因子。
另一方面，表 \ref{tab:fold_bin_degeneracy_spectrum_m6_m12} 给出 \(d_{\max}(10)=9\)，而定理 \ref{thm:fold-bin-min-degeneracy-fib-audited} 给出 \(d_{\min}(12)=8\)，故 \(G_{10}^{\mathrm{bin}}\) 与 \(G_{12}^{\mathrm{bin}}\) 都不可解。
\end{theorem}
\begin{proof}
由命题 \ref{prop:fold-bin-gauge-decomposition}，
\[
G_m^{\mathrm{bin}}
\cong
\prod_{w\in X_m}\mathfrak S_{d_w}.
\]
对有限直积群有
\[
Z\!\left(\prod_{w\in X_m}\mathfrak S_{d_w}\right)
\cong
\prod_{w\in X_m}Z(\mathfrak S_{d_w}).
\]
而
\[
Z(\mathfrak S_d)\cong
\begin{cases}
\ZZ/2\ZZ,& d=2,\\
1,& d\neq 2,
\end{cases}
\]
故立得 \(Z(G_m^{\mathrm{bin}})\cong (\ZZ/2\ZZ)^{B_m}\)。

有限直积群可解当且仅当每个因子都可解，而 \(\mathfrak S_d\) 可解当且仅当 \(d\le 4\)。
因此
\[
G_m^{\mathrm{bin}}\ \text{可解}
\iff
d_{\max}(m)\le 4.
\]
在可解情形下，导出长度等于各因子导出长度的最大值。
再用
\[
\ell_{\mathrm{der}}(\mathfrak S_1)=0,\qquad
\ell_{\mathrm{der}}(\mathfrak S_2)=1,\qquad
\ell_{\mathrm{der}}(\mathfrak S_3)=2,\qquad
\ell_{\mathrm{der}}(\mathfrak S_4)=3,
\]
便得到所示分段公式。

\(m=6\) 的结论由表 \ref{tab:fold_bin_degeneracy_spectrum_m6_m12} 的直方图 \(2:8,3:4,4:9\) 直接给出，而 boundary 中心子群为规范直和因子正是定理 \ref{thm:xi-window6-center-exact-splitting-bdry-cyclic} 的内容。
对 \(m=10\)，表中存在 \(d=9\) 纤维，故含不可解因子 \(\mathfrak S_9\)。
对 \(m=12\)，由定理 \ref{thm:fold-bin-min-degeneracy-fib-audited} 的 \(d_{\min}(12)=8\) 可知每条纤维都满足 \(d_w\ge 8\)，从而同样不可解。证毕。
\end{proof}

\endinput
